% Source r6 scalar-mass versions and examples.
\section{Excluding finite normalised negative mass}\label{sec:mass}

Theorem~\ref{res:bounded} uses a uniform bound on the negative part.  A different sufficient hypothesis, on the total normalised negative mass,
gives the same endpoint.  Its proof is multiplicative rather than arithmetic:
it bounds $C$ by a product, then uses integer discreteness.

\begin{theorem}[finite normalised negative mass]\label{res:mass}
Let \(C_n\in\N_{>0}\) and \(E_n\in\Z\) satisfy
\[
 C_{n+1}=C_n-E_n,\qquad
 \forall K\ \exists N\ \forall n\ge N,\quad K|E_n|<C_n .
\]
If
\[
 \sum_{n=0}^{\infty}\frac{(-E_n)_+}{C_n}<\infty ,
\]
then \(E_n=0\) eventually.  For an exact reciprocal-tail orbit this implies
\(a_{n+1}=a_n^2-a_n+1\) eventually.
\end{theorem}

\begin{proof}
Put \(\delta_n=(-E_n)_+/C_n\).  The update gives
\(C_{n+1}\le C_n(1+\delta_n)\), so
\[
 C_N\le C_0\prod_{n<N}(1+\delta_n).
\]
The product is bounded because \(\sum\delta_n\) converges.  Choose an integer
\(M\) above that bound and apply normalised vanishing with \(K=M\).  Then
\(M|E_n|<C_n<M\) for all sufficiently large \(n\), forcing the integer
\(E_n\) to be zero.  The final recurrence is Theorem~\ref{res:eventual}, whose
hypothesis $C_{n+1}\ne0$ holds because $C_n>0$ throughout.
\end{proof}

\noindent The same product already caps $C_n$, so a strict integer rise costs a
definite relative mass $\ge 1/K$ and only finitely many rises can occur.  After
the last rise the tail is nonincreasing, hence eventually constant.  Normalised
vanishing is not required for that scalar conclusion.

\begin{lemma}[scalar finite mass]\label{res:massscalar}
Let \(C_n\in\N_{>0}\) and \(E_n\in\Z\) satisfy \(C_{n+1}=C_n-E_n\).  If
\(\sum_n(-E_n)_+/C_n<\infty\), then \(E_n=0\) eventually.  For an exact
reciprocal-tail orbit this implies \(a_{n+1}=a_n^2-a_n+1\) eventually.
\end{lemma}

\noindent Formalised as the
\lword{Erdos243/SparseResetRecovery.lean}{551}{bounded_of_summable_negativeRelativeMass}{product
bound on \(C\)}, the
\lword{Erdos243/SparseResetRecovery.lean}{612}{eventually_zero_of_summable_negativeRelativeMass_scalar}{scalar
vanishing without normalised vanishing}, and the
\lword{Erdos243/SparseResetRecovery.lean}{691}{sylvesterNext_eventually_of_summable_negativeRelativeMass_scalar}{Sylvester
consumer without vanishing}.  Theorem~\ref{res:mass} remains the registered
form that still lists vanishing, matching the earlier Lean declarations.

\noindent The product argument is the
\lword{Erdos243/SparseResetRecovery.lean}{90}{eventually_zero_of_summable_relativeGrowth}{summable
relative growth}; its specialisation to negative mass is the
\lword{Erdos243/SparseResetRecovery.lean}{155}{eventually_zero_of_summable_negativeRelativeMass}{vanishing
from finite negative mass}, and the statement giving the recurrence directly is
the
\lword{Erdos243/SparseResetRecovery.lean}{175}{sylvesterNext_eventually_of_summable_negativeRelativeMass}{Sylvester
recurrence from summable negative mass}.  This is a checked implication about
the state system, not a derivation of normalised vanishing from
Problem~\#243.  Normalised vanishing is still an input of
Theorem~\ref{res:mass}: it is what converts the bound on $C$ into a bound on
$|E_n|$.  Lemma~\ref{res:massscalar} does not use it.  Nothing here derives
summability from the growth condition of Problem~\ref{res:problem}.

The transfer of the previous section applies here as well.  On the canonical
orbit of Corollary~3 of~\cite[p.~9]{koizumi2025} the recurrence
$C_{n+1}=C_n-E_n$ holds with $C_n>0$ by Lemma~4
~\cite[p.~11]{koizumi2025}, and normalised vanishing is the vanishing of
the gap sequence, so Theorem~\ref{res:mass} is a conditional theorem about
Problem~\ref{res:problem} whose sole remaining hypothesis is summability of
$\sum_n(-E_n)_+/C_n$.  That hypothesis is not supplied there, and is the
content of Problem~\ref{res:masshyp} below.

\begin{example}[the product bound at a geometric rate]\label{ex:mass}
Suppose $C_0=100$ and $(-E_n)_+/C_n\le2^{-n-1}$ for every $n$, so that the
normalised negative mass is at most $1$; the constraint at $n=0$ still permits
$E_0$ as negative as $-50$.  Since $1+x\le e^{x}$,
\[
 C_N\le100\prod_{n<N}\bigl(1+2^{-n-1}\bigr)\le100e<272
\]
for every $N$.  Normalised vanishing at $K=272$ then gives
$272|E_n|<C_n<272$ for all large $n$, so $|E_n|<1$ and $E_n=0$ there.  The
constant $272$ comes from the total mass and not from any single magnitude,
which is why the hypothesis of Theorem~\ref{res:mass} is summability of
$(-E_n)_+/C_n$ and not a bound on it.
\end{example}
